% 本文件由 LaTeX 排版 Agent 根据有效 translation.json 自动生成。
% author=Vincent of Lérins (d. c. 450); work_id=3506; title=备忘录

\chapter{\WorkTitle{备忘录}}

% structure: p0001 source=h1 target=omitted reason=page-title-already-rendered-by-parent

\Emph{为了天主教信仰的古老性和大公性，反对所有异端亵渎性新奇的著作。}\par

% structure: p0003 source=h2 target=section reason=relative-heading-rank; base=h2

\section{第一章}

\Emph{本文目的。}\par

[1.] \Emph{我，\Person{Peregrinus}}，是\Person{天主}所有仆人中最小的，想起圣经的劝诫：\Quote{你当询问你的父亲，他们必告诉你；你的长者，他们必对你说明。}\BibleRef{申 32:7} 又说：\Quote{你要侧耳倾听智者的言辞，}\BibleRef{箴 22:17} 再次：\Quote{我儿，不要忘记这些训诲，愿你的心谨守我的诫命。}\BibleRef{箴 3:1} 我，\Person{Peregrinus}，因这些劝诫，深信在主助佑下，将我从圣教父们那里真实领受的事记录下来，并非毫无益处，尤其对我微薄的能力而言，实为必要；这样我便能通过不断阅读，弥补记忆的不足。\par

[2.] 激励我的不仅是对劳苦果实的期盼，还有对时间及地方适宜性的考虑：\par

对时间的考虑——因为眼见时间攫取一切人间事物，我们亦当从中夺取一些能有益于永生的事物，尤其当天主审判临近的可怕期待迫切要求我们更专务于信仰，而新异端者的诡诈狡猾更需要特别的警醒和关注。\par

地方适宜性也激励着我：我避开城市的喧嚣和人群，住在一座隐修院的静谧中，该院位于偏远的庄园；在那里，我能无分心地遵从圣咏作者的劝诫：\Quote{你们要静默，并承认我是\Person{天主}。}\par

此外，这也与我的生活目的相符：我曾一度卷入世俗战争中多变而可悲的风暴，如今终于在基督的保佑下，将锚抛入宗教的港湾——这港湾对所有人永远最为安全——以便在那里摆脱虚荣和骄傲的狂风，以基督徒谦卑的祭献平息天主的义怒，从而不仅逃脱今生的船难，也逃避来世的火焰。\par

[3.] 但现在，我要以主的名义着手实现我的目标：即，以叙述者的忠诚，而非作者的妄自尊大，记录我们的祖先所传授并托付给我们的事；但在写作中遵循这一原则：绝不完全触及所有可说的事，只涉及必要之事；亦不以华丽精确的文风，而以简单通俗的语言，使大部分内容看似暗示而非详述。让那些对自己的能力有信心或因责任感而受感动的人去追求优雅和精确。对我来说，为自己预备一份\Emph{备忘录}（或提醒录）就足够了，以帮助我的记忆，更确切地说，防止我的遗忘；然而，我将努力在主助佑下，通过回忆所学的知识，日复一日地逐渐修改和完善这份备忘录。我事先提及这一点，倘若我写的东西偶然遗失，落到圣人们手中，他们看到尚有改进和完善的承诺，便不会匆忙指责其中任何内容。\par

% structure: p0011 source=h2 target=section reason=relative-heading-rank; base=h2

\section{第二章}

\Emph{辨明天主教信仰真理与异端者邪恶虚假的一般规则。}\par

[4.] 因此，我常常殷切而专注地询问许多以圣德和学识著称的人，如何通过某种可靠、可谓普遍的规则，辨明天主教信仰的真理与异端者的邪恶虚假；我几乎总是得到这样的答复：无论我或任何人，若要识别异端者兴起时的欺诈、避开他们的陷阱，并在天主教信仰中保持健全和完整，就必须在主助佑下，以两种方式强化自己的信念：首先，凭天主法律的权威；其次，凭天主教会传承的权威。\par

[5.] 但这里或许有人会问：既然圣经正典是完备的，本身已足够一切，甚至绰绰有余，为何还需要教会的解释权威呢？理由如下：由于圣经的深奥，并非所有人都以同一意义接受它，而是各有理解；因此，它似乎有多少解释者，就能有多少种解释。因为\Person{诺瓦替安}以一种方式解释，\Person{撒伯流}以另一种，\Person{多纳图}以一种，\Person{亚略}、\Person{欧诺米}、\Person{马其顿尼}以另一种，\Person{福提努}、\Person{阿波利纳利}、\Person{普里西利安}以另一种，\Person{约维尼安}、\Person{佩拉纠}、\Person{塞莱斯提}以另一种，最后\Person{聂斯托利}又是一种。因此，面对如此众多错误的复杂情况，非常有必要按照教会和天主教解释的标准，来确立正确理解先知和宗徒的规则。\par

[6.] 此外，在教会本身内，必须尽一切可能持守那处处、时常、众所信仰的信仰。因为这才是真正且严格意义上的\Quote{公教}信仰——正如其名称和事物本身所表明的——它普遍包含一切。我们将遵循这一规则，如果我们追随大公性、古老性、共识性。如果我们承认那普世教会所承认的一个信仰为真，我们就追随了大公性；如果我们决不背离我们圣洁的祖先和教父们所明显持有的解释，我们就追随了古老性；同样地，如果我们坚持古老性本身中所有、或至少几乎所有司铎和圣师的一致界定和决定，我们就追随了共识性。\par

% structure: p0016 source=h2 target=section reason=relative-heading-rank; base=h2

\section{第三章}

\Emph{如果有一个人或多人不同意其余人，应当怎么做。}\par

[7.] 那么，若教会的微小部分已自普世信仰的共融中割裂，一位天主教基督徒该当如何？无疑，岂非宁可选择全体肢体的健全，而胜过染疫腐坏的肢体？若某种新奇的传染试图感染，不是教会中微不足道的一部分，而是全体，那又如何？那时，他当留意持守古老传统，因为古老传统今日决不可能被任何新奇的欺骗所诱惑。\par

[8.] 但是，若在古老传统本身之中发现有两三个人的错误，或至少是一个城市甚至一个省份的错误，那又如何？那时，他当务必优先遵循古代大公会议的法令（若有的话），而非少数人的鲁莽和无知。但是，若出现某种错误，而关于该错误并未发现有任何这样的法令，那他就当汇集、咨询并查问古代教父的意见，即那些虽生活于不同时代和地点，却仍在同一大公教会的共融和信仰中持续，并被公认为权威者；并且，无论他查明什么是被他们全体一致、公开、频繁、持久地持有、书写、教导的，而不只一两个人，他就当明白，自己也当毫无怀疑或犹豫地相信那一点。\par

% structure: p0020 source=h2 target=section reason=relative-heading-rank; base=h2

\section{第四章}

\Emph{因引入新奇教义所产生的邪恶，以多纳徒派和亚略派的实例说明。}\par

[9.] 为使我们的论述更易理解，须以个别实例加以说明，并稍作详细阐述，以免因过于简短而忽略重要之事。\par

在\Person{多纳徒}的时代——其追随者称为多纳徒派——当\Place{非洲}大量人鲁莽地冲入他们自己的疯狂错误，并忘记自己的名号、信仰和宣认，宁愿以一个人的亵渎妄为取代基督的教会时，那时全\Place{非洲}唯有那些厌弃这亵渎的分裂、继续与普世教会共融的人，才安全地处于大公信仰的圣所之内；他们为后人留下了杰出的榜样，说明将来当如何以及何等明智地将全体的健全置于一人或至多少数人的疯狂之上。\par

[10.] 同样，当亚略的毒害感染了教会中并非微不足道的一部分，而是几乎全世界，以至于一种盲目似乎落在了几乎所有的拉丁语主教身上——他们部分被武力、部分被诡计所欺瞒——使他们无法在如此混乱中看到最适宜的做法时，那时，凡是真实热爱并朝拜基督的人，宁愿选择古老的信仰而舍弃新奇的谬误，便逃脱了那致命的感染。\par

[11.] 那时代的危险充分表明，引入新奇教义会带来多大的灾难。因为那时，不仅微不足道的利益，而且其他极为重要的事物都被颠覆了。因为不仅姻亲、关系、友谊、家族，而且还有城市、人民、省份、民族，最后整个罗马帝国，都被动摇根基而毁灭。因为当这同一亵渎的亚略新奇教义，如同战争女神或复仇女神，先是俘获了皇帝，然后又使宫中所有主要人物服从新律法时，从那时起，它便从未停止将一切卷入混乱，搅乱所有事物，公共的和私人的，神圣的和亵渎的，不顾良善与真实，反而如同拥有权威地位，随意击打任何人。于是妻子被玷污，寡妇被蹂躏，童贞女被亵渎，隐修院被拆毁，神职人员被驱逐，下级神职人员被鞭打，司铎被流放，监狱、牢房、矿场充满了圣徒，他们中的大多数被禁止进入城市，被逐出家园，流浪在沙漠和洞穴中，在岩石和野兽出没之地，赤身露体、饥饿、干渴而憔悴耗尽。凡此种种，除了因为人的迷信被引入以取代天上的教义，因邪恶的新奇颠覆了确立已久的古老传统，因前代的制度被藐视，因祖先的法令被废除，因先辈的决定被撕碎，以及亵渎新奇的贪欲拒绝将自己限制在神圣无玷古老传统的最贞洁界限之内，难道还有其他原因吗？\par

% structure: p0026 source=h2 target=section reason=relative-heading-rank; base=h2

\section{第五章}

\Emph{殉道者给我们树立的榜样：任何强权都不能阻止他们捍卫先辈的信德。}\par

[12.] 但或许，我们出于对新奇的憎恨和对古老传统的热忱而捏造了这些指控。若有人愿意听信这种暗示，至少当相信真福\Person{安博}。他在写给\Person{格拉提安}皇帝的著作的第二卷中，哀叹当时的严酷，说：\Quote{如今已足够，全能的天主！我们以自己的毁灭、自己的血，赔补了宣信者的被杀、司铎的流放和如此极端不敬的邪恶。无疑，已确知，那些违犯信德的人不能安然无恙地存留。}\par

在同一著作的第三卷中，他说：\Quote{让我们遵守前人的命令，不以粗鲁的鲁莽越过我们从他们继承的界标。那密封的预言书，没有长者、没有权能者、没有天使、没有总领天使敢于打开。唯独基督拥有解释它的特权。\BibleRef{默 5:1}我们中谁敢于开封那被宣信者封印、并已由众多殉道者的致命所祝圣的司铎之书？那些被迫力开封的人事后又重新封印，谴责加于他们的欺骗；而那些不敢动手开封的人则证明自己是宣信者和殉道者。我们怎能否认那些我们宣扬其胜利之人的信德？}\par

[13.] 我们真实地宣讲，可敬的安波罗削，我们宣讲，并赞赏与钦佩。因为谁如此疯狂，虽不能赶上，却至少不热切渴望跟随那些人？没有任何力量——无论是威胁、谄媚、生命、死亡、皇宫、御林军、皇帝、帝国本身、人、魔鬼——能阻止他们捍卫祖先的信德。我再说，主因他们坚守宗教古训，使他们堪受如此巨大的赏报，借着他们的手，主重建了被毁的教会，使灵性死亡的人民获得新生，将夺去的司铎冠冕重新戴在他们头上，并以信德的泪水（那是天主在主教们心中开启的源泉）洗去了那些可憎的——我不说文字，而是\Emph{涂改}（\OriginalTerm{不是文字，而是涂改}{non literas, sed lituras}）——新奇的邪恶。最后，当几乎整个世界被意料之外的异端无情风暴压倒时，主借着他们，使世界从新奇的谬误回归古老的信德，从新奇的疯狂回归古训的健全，从新奇的盲目回归最初的光明。\par

[14.] 然而，在这些宣信者所展现的（我们可称之为）神圣勇毅中，我们必须特别注意，他们当时诉诸古代教会而进行的辩护，并非为部分，而是为全体辩护。因为如此卓越的人士，不应竭尽全力维护一两个人的模糊而矛盾的观点，也不应为某个微不足道省份的轻率联盟而奋斗；相反，他们坚守神圣教会全体司祭阶层的法令与定义，作为宗徒与公教真理的继承者，他们宁愿交出自己，也不愿出卖大公性与古老性的信德。因此，他们被认为配得如此伟大的光荣，不仅被视为宣信者，而且正确而应得地被尊为宣信者中的佼佼者。\par

% structure: p0032 source=h2 target=section reason=relative-heading-rank; base=h2

\section{第六章}

\Emph{教宗斯德望抵制重洗的榜样。}\par

[15.] 因此，这些蒙福之人的榜样是伟大的，这榜样显然神圣，值得每一位真正的公教信徒不断回忆和默想。他们如同七灯台，闪耀着圣神的七重光明，向后世指明：从此以后，亵渎的新奇之狂妄，在种种谬误的喧嚣中，如何能被神圣古训的权威所粉碎。\par

这并非什么新鲜事？因为在教会中，一向如此：一个人越受宗教影响，就越急于反对革新。这样的例子数不胜数；但为简洁起见，我们只取一个，而且优先取自宗座，以便对每个人都显而易见：蒙福的宗徒们的蒙福继承人，是何等有力地、何等热切地、何等恳切地捍卫他们曾一次领受的宗教的完整。\par

[16.] 曾经，可敬的迦太基主教阿格里皮努斯持守一项教义——他是第一个持守的人——即应当重行圣洗，这违反神圣的法则，违反普世教会的规则，违反我们祖先的习俗与制度。这项革新带来了如此巨大的邪恶，以致不仅为各种异端者提供了亵渎的范例，也成了某些公教信徒犯错的机会。\par

当众人抗议这新奇之事，各地司祭阶层按各自的热诚纷纷反对时，蒙福的教宗斯德望（可敬的宗座主教）与他的同僚们一起，却自己首当其冲地抵制了它。我毫不怀疑，他认为自己既在地位权威上超越众人，也应在信德的奉献上超越众人。最终，在当年致非洲的信函中，他定下这条规则：\Quote{不得有任何革新——只应持守所传授的。}因为这位圣善而明智的人深知，真正的虔诚不接受任何其他规则，只接受：凡从祖先忠信领受的一切，都应忠信地传给子女；我们的责任不是引领信仰随我们的意愿，而是跟随信仰所引领的；基督徒的谦逊与庄重在于：不将我们自己的信念或习规传给后代，而应保存和持守我们从先辈领受的一切。那么，这整个事件的结果是什么？不过是惯常的结果：古训被保留，新奇被拒绝。\par

[17.] 但也许，当时革新之事业缺乏支持。恰恰相反，它拥有如此强大的才华、如此丰富的口才、如此众多的拥护者、如此像真理的外表、以及圣经（只是以新奇而乖戾的方式解释）如此有力的支持，以致在我看来，那整个阴谋若非因这场特别动荡的唯一原因——即所如此从事、如此辩护、如此赞扬之事的本身的新奇性——使之下滑，它本不可能被击败。最终，同上那个非洲公会议或其法令，在天主眷顾下，产生了什么结果？毫无结果。整件事如同一个梦、一个寓言、一件无足轻重之事，被废除、取消并践踏于脚下。\par

[18.] 啊，何等奇妙的转变！这同一教义的作者被视为公教信徒，而追随者却成了异端者；教师被赦免，门徒被定罪；著作的作者将是天国之子，而辩护者将在地狱中有份。因为谁如此疯狂，竟怀疑所有圣洁主教与殉道者中的蒙福之光——西彼廉，连同他的同僚们——将与基督一同为王？而谁又如此亵渎，竟否认多纳徒派以及其他那些害虫（他们夸耀那公会议的权威来支持重行圣洗），将与魔鬼一同被投入永火？\par

% structure: p0040 source=h2 target=section reason=relative-heading-rank; base=h2

\section{第七章}

\Emph{异端者如何狡猾地引用古代作者中的隐晦段落来支持自己的新奇说法。}\par

[19.] 这项谴责似乎是天主上智颁布的，尤其针对那些人的欺诈：他们设法以异于其本名的名义来粉饰异端，常常抓住某位古代作者表达不甚清晰的作品，由于自己教义本身的晦涩，这些作品看起来与其相符，从而获得他们既非首创亦非唯一持有此说的名声。依我之见，他们的这种恶行加倍可憎：第一，因为他们不惧怕邀请他人饮下异端的毒药；第二，因为他们以渎神的气息，仿佛扇动余烬成火焰，吹向每位圣人的记忆，并通过重新提起本应缄默埋葬的事而散布恶名，如此踏上了其父含的后尘：含不仅不遮盖可敬诺厄的裸体，反而告诉其他人以供嗤笑，如此严重地违反孝道义务，以至于连他的后代都与他一同承受他所招致的诅咒，与他那些蒙福的兄弟大相径庭：那些兄弟既不愿以眼玷污敬仰父亲的裸体，也不容许他人这样做，而是如经上所说，倒退着走去遮盖他，即他们既不赞同也不暴露圣人的过错，为此他们自己和后裔都得到祝福的赏报。\BibleRef{创 9:22}\par

[20.] 但回到当前之事：我们必须非常畏惧败坏信仰和玷污宗教的罪行，不仅教会的纪律，而且权威的谴责都阻止我们犯此罪。因为众所周知，蒙福的宗徒保禄是何等严厉、何等剧烈、何等强烈地抨击某些人，他们以惊人的轻率，\Quote{这么快就离开了那召你们进入基督恩宠的，去归向另一个福音；那并不是另一个福音}\BibleRef{迦 1:6}，\Quote{他们为自己的私欲增添了师傅，掩耳不听真理，偏向无稽的传说}\BibleRef{弟后 4:3}，\Quote{他们因废弃了起初的信德而受惩罚}\BibleRef{弟前 5:12}，他们被那些人欺骗，关于那些人，同一宗徒致罗马基督徒写道：\Quote{弟兄们，我劝你们，要提防那些造成分裂和绊脚石、反对你们所学的道理的人，你们要避开他们。因为这样的人不服事我们的主基督，只服事自己的肚腹，用甜言蜜语欺骗老实人的心}\BibleRef{罗 16:17}，\Quote{他们潜入人家，迷惑那些被各种私欲牵引、被罪担累的软弱妇女，她们常常学习，终久不能明白真道}\BibleRef{弟后 3:6}，\Quote{空谈家与欺骗者，他们败坏整个家庭，教导不当教的事，为可耻的利益}\BibleRef{铎 1:10}，\Quote{他们心思败坏，在信仰上可废弃}\BibleRef{弟后 3:8}，\Quote{傲慢无知，专好辩论和言词上的争执，缺乏真理，以虔敬为得利的门路}\BibleRef{弟前 6:4}，\Quote{同时他们又学习懒惰，挨家闲游，不但懒惰，而且说长道短，好管闲事，说不该说的话}\BibleRef{弟前 5:13}，\Quote{他们丢弃了良心，就在信仰上翻了船}\BibleRef{弟前 1:19}，\Quote{他们亵渎的空谈日益增加不虔敬，他们的话如同毒疮蔓延}\BibleRef{弟后 2:16}。关于他们，也写得很好：\Quote{但他们不能再前进，因为他们的愚昧将暴露在众人面前，如同那两个人一样}\BibleRef{弟后 3:9}。\par

% structure: p0044 source=h2 target=section reason=relative-heading-rank; base=h2

\section{第八章}

\Emph{圣保禄宗徒话语的解释，\BibleRef{迦 1:8}}\par

[21.] 因此，当这类人游历各省各城，带着他们可贩卖的错谬，来到迦拉达，迦拉达人听了他们之后，厌恶真理，呕吐出宗徒及大公教会的教义这玛纳，喜爱异端的新奇垃圾时，宗徒运用他职务的权威，以极严厉的措辞宣布判决：\Quote{即使是我们，或是从天上来的天使，若向你们宣讲我们所传的福音以外的福音，该受诅咒}\BibleRef{迦 1:8}。\par

[22.] 他为什么说\Quote{即使是我们}？为什么不说\Quote{即使是我}？他的意思是：\Quote{即使是伯多禄、安德肋、若望，总之，即使是整个宗徒团体，向你们宣讲我们所传的以外的福音，也该受诅咒。}何等可怕的严厉！他既不宽恕自己，也不宽恕他的同宗徒，只为保存最初传授的信仰毫不变更。不仅如此，他还继续道：\Quote{即使是天上的一位天使，若向你们宣讲我们所传的福音以外的福音，也该受诅咒。}为保存一次传授的信仰，仅仅提到人是不够的，还必须包括天使。\Quote{即使是我们，}他说，\Quote{或是从天上来的天使。}这并不是说天上的圣天使如今可能犯罪。而是他的意思是：即使发生那不可能发生的事——如果有人，无论他是谁，企图改变一次传授的信仰，就该受诅咒。\par

[23.] 但也许他起初这样说是不加思索的，发泄了人的冲动，而不是在神圣引导下表达自己。绝非如此。他接着自己所说的，并以强烈反复的热忱敦促说：\Quote{正如我们先前说过的，现在我再说，若有人给你们传讲另一个福音，与你们所领受的不同，当受咒诅。} 他没有说：\Quote{若有人给你们传讲另一信息，与你们所领受的不同，愿他受祝福、被称赞、受欢迎，}——不；而是\Quote{当受咒诅，}即\OriginalTerm{绝罚}{anathema}，被隔离、被分开、被排除，以免一只羊的致命感染通过其有毒的混杂污染基督的无辜羊群。\par

% structure: p0049 source=h2 target=section reason=relative-heading-rank; base=h2

\section{第九章}

\Emph{他对迦拉达人的警告是对所有人的警告。}\par

[24.] 但也许这警告只是针对迦拉达人。就算是这样；那么在同一封书信中随后的其他劝诫也只是针对迦拉达人的，例如：\Quote{如果我们活在圣神内，就让我们也随从圣神而行；不要贪图虚荣，彼此惹气，互相嫉妒。}等等；\BibleRef{迦 5:25} 如果这种看法是荒谬的，而诫命是同样针对所有人的，那么由此可知，正如这些关于伦理的诫命同样针对所有人，那些关于信仰的警告也同样针对所有人；并且正如所有人都不可彼此惹气或互相嫉妒，同样，所有人都不可接受任何不同于天主教会到处宣讲的福音。\par

[25.] 或者，那对宣讲不同于所传福音者的绝罚是针对当时，而非现今。那么，\Quote{随从圣神而行，就不会满足肉体的私欲}（\BibleRef{迦 5:16}）的劝诫也是针对当时，而非现今。但如果相信这点既是不敬也是有害的，那么必然得出，正如这些诫命为所有时代所遵守，那些禁止改动信仰的警告也是为所有时代而发的警告。因此，向天主教基督徒宣讲任何不同于他们所领受的道理，过去从未合法，现在也从未合法，将来也绝不合法；而绝罚那些宣讲任何不同于一次所领受的人，过去常是责任，现在常是责任，将来也常是责任。\par

[26.] 既然如此，有谁如此胆大妄为，竟宣讲不同于教会所传的道理，或有谁如此反复无常，竟接受不同于他从教会所领受的道理呢？那蒙拣选的器皿，那外邦人的导师，那宗徒的号角，那奉召向全地宣讲的传道者，那被提到第三层天上的人（\BibleRef{格后 12:2}），在他的书信中向所有人、时常、各处一再呼喊：\Quote{若有人宣讲任何新道理，当受咒诅。} 另一方面，一群短暂、垂死的青蛙、跳蚤和苍蝇，诸如\OriginalTerm{白拉奇派}{Pelagians}，却起来反对，并向着天主教徒说：听从我们的话，跟随我们的引导，接受我们的解释，谴责你们曾经持有的，持有你们曾经谴责的，抛弃古老的信仰、你们祖先的训导、你们先辈留给你们的托付，然后接受——什么？我战栗不敢说出：因为它充满如此傲慢和自负，以致在我看来，不仅肯定它，甚至驳斥它，都不能在某种程度上没有罪过。\par

% structure: p0054 source=h2 target=section reason=relative-heading-rank; base=h2

\section{第十章}

\Emph{为何天主允许卓越人物成为教会中新奇的倡导者。}\par

[27.] 但有人会问：既然如此，为何教会中某些优秀且有地位的人，时常被天主允许向天主教徒宣讲新奇道理呢？这当然是一个恰当的问题，必须非常仔细和充分地处理，但回答时不应依赖自己的能力，而应凭借神圣法律的权威，并诉诸教会的决定。\par

那么，让我们聆听圣梅瑟，让他教导我们，为何学识渊博、且因其知识甚至被宗徒称为先知的人，有时被允许提出新奇道理——旧约常以寓意的方式称之为\Quote{异神}，因为异端者对其观念所怀的敬意，与外邦人对其神祇所怀的敬意相同。\par

[28.] 蒙福的梅瑟遂在\WorkTitle{申命纪}中写道：\BibleRef{申 13:1-3} \BibleQuote{如果你们中间出现了一位先知或作梦的人，}即在教会中担任圣师职的人，被他的门徒或听众认为是由启示而教导的；那么——接下来呢？\BibleQuote{他给你一个征兆或奇迹，并且他所说的征兆或奇迹实现了，}——他是指某位杰出的圣师，其学识如此渊博，以致他的追随者相信他不仅知晓人间之事，而且还能预知超人之事，正如他们的门徒通常夸耀的那样，有\Person{瓦伦提诺}、\Person{多纳图}、\Person{福提诺}、\Person{阿波利纳里}以及那一类人！接着呢？\BibleQuote{并且对你说：我们去跟随别的神，就是你们所不认识的神，并事奉他们。}那些别的神是什么？不就是你们不认识的外邦谬误吗？也就是前所未闻的新异端？\BibleQuote{我们事奉他们；}也就是说，\Quote{我们相信他们，跟随他们。}最后呢？\BibleQuote{你不要听那先知或作梦之人的话。}请问，天主为什么不准人教导他禁止人听的事呢？\BibleQuote{因为上主你们的天主是在试探你们，要知道你们是否全心、全灵爱慕祂。}天主眷顾有时允许某些教会圣师宣讲新道理的原因比白天还清楚——\BibleQuote{这是上主你们的天主在试探你们}；祂说。这确实是一个重大的试探：当你相信一个人是先知，是先知的门徒，是真理的教师和捍卫者，以极大的敬爱与爱慕将他怀抱在胸前，而这样的人突然之间暗中诡秘地引入有害的谬误，这些谬误你既因被他先前权威的声望所束缚而无法迅速识破，又因对旧有恩师的感情阻碍而不能轻易认为应该加以谴责。\par

% structure: p0059 source=h2 target=section reason=relative-heading-rank; base=h2

\section{第十一章}

\Emph{来自教会历史的例证，证实梅瑟的话——\Person{聂斯多略}、\Person{福提诺}、\Person{阿波利纳里}。}\par

[29.] 或许有人会要求我们用教会历史的例证来阐明圣梅瑟的话。这个要求是合理的，也不会久等而得满足。\par

[29.] 首先举一个最近且很明显的事例：我们想，那是怎样的一种试探，不久前教会所经历的，当那个不幸的\Person{聂斯多略}，突然从羊变成狼，开始蹂躏基督的羊群，而他所吞食的羊群中，有很大一部分仍相信他是羊，因此更容易受到他的攻击？因为谁会轻易相信他错了呢？他是由皇帝的高选所拣选，并深受司祭团的尊崇。谁又会轻易相信他错了呢？他深受圣洁弟兄们的爱戴，在民众中极受欢迎，每日公开讲解圣经，并驳斥犹太人和外邦人的有害谬误。谁能不相信他的教导是正统的，他的讲道是正统的，他的信仰是正统的呢？他为了给自己的异端开道，还热切地抨击所有异端的亵渎。但这正是梅瑟所说的：\BibleRef{申 13:3} \BibleQuote{上主你们的天主试探你们，要知道你们是否爱祂。}\par

[30.] 撇开\Person{聂斯多略}（他身上总是令人赞赏多于获益，炫耀多于真实，他凭天赋能力而非天主恩宠在一般人眼中得势一时），我们转而谈论那些学识渊博、勤勉过人的人，他们曾给公教徒带来不小的考验。例如，在\Place{潘诺尼亚}的\Person{福提诺}，据我们父辈的记忆，他曾在\Place{西尔米乌姆}教会成为一次考验。他在那里受到普遍赞许而晋升司铎，并以公教徒身份履职一段时间后，突然，如同梅瑟所说的那邪恶的先知或作梦的人一样，他开始劝诱托付给他照管的百姓跟随\BibleRef{申 13:2} \BibleQuote{外邦神祇}，即他们先前不认识的外邦谬误。但这并无异常：问题的祸害在于，他为了实现如此重大的邪恶，利用了非同寻常的辅助。因为他天赋甚高，口才雄辩，学识渊博，用两种语言辩论和著述丰富而有力，他遗留下来的著作，一部分用希腊文写成，一部分用拉丁文写成，可资证明。幸而托付给他的基督的羊群，为公教信德保持警醒和谨慎，迅速将目光转向梅瑟的告诫之言，虽然钦佩他们先知和牧者的口才，却未忽视考验。因为从那时起，他们开始躲避他如狼，而先前他们曾跟随他如群羊之首。\par

[31.] 我们不仅从\Person{福提诺}的实例中认识到这种考验对教会的危险，同时也得到提醒，在护卫信德时需要加倍勤勉。\Person{阿波利纳里}也发出了类似的警告。因为他使他的门徒们产生了极大的争议和严重的困惑，教会的权威将他们拉向一边，他们老师的势力却拉向另一边；因此他们在两者之间摇摆不定，不知所措。\par

但或许他是一位品性毫无分量的人呢？恰恰相反，他是如此卓越和受人敬重，以至于他所言之事极易被人接受。因为还有什么能超过他的敏锐、他的机智、他的学识呢？他在许多著作中摧毁了多少异端！他驳斥了多少敌对信仰的错误！证据就是那部最崇高、最广博的著作，不下三十卷，以大量的论证击退了\Person{波斐利}的疯狂诽谤。要列举他所有的著作将耗费很长时间，这些著作肯定会使他位列教会建造者中的翘楚，若非那种异端好奇心的亵渎欲望引导他发明我不知道的某种新奇事物，那种新奇事物如同麻风病般的传染玷污了他所有的辛劳，并导致他的教导被宣判为教会的考验而非教会的建造。\par

% structure: p0066 source=h2 target=section reason=relative-heading-rank; base=h2

\section{第十二章}

\Emph{关于\Person{斐提努斯}、\Person{阿波利纳里斯}和\Person{聂斯多略}之错误的更详细叙述}\par

[32.] 这里，或许有人会问及上述异端的若干叙述；即\Person{聂斯多略}、\Person{阿波利纳里斯}和\Person{斐提努斯}的异端。确实，这并非当前议题；因为我们的目的不是详述个别人的错误，而是提供少数例子，在其中可以清楚明显地看到\Person{梅瑟}言语的适用性；也就是说，如果有时某位教会大师，他本人也是解释先知奥秘的先知，试图在教会中引入某种新教义，天主上智允许此事发生是为了考验我们。因此，作为插叙，简要说明上述异端者\Person{斐提努斯}、\Person{阿波利纳里斯}、\Person{聂斯多略}的意见将是有益的。\par

[33.] \Person{斐提努斯}的异端如下：他说天主是独一唯一，如同犹太人所认为的。他否认天主圣三的完满，不相信有圣言的位格或圣神的位格。他断言基督只是一个普通人，其本源来自玛利亚。因此，他极其固执地主张我们只应崇拜圣父的位格，而只应以人的身份尊敬基督。这就是\Person{斐提努斯}的教义。\par

[34.] \Person{阿波利纳里斯}假装在关于天主圣三的统一性上赞同教会（尽管连这一点他也不完全持纯正信念），但在关于主的降生成人上则公开亵渎。因为他说我们救主的肉身或者完全没有人的灵魂，或者至少没有理性的灵魂。此外，他说主的同一肉身并非取自童贞玛利亚的血肉，而是从天上降入贞女；并且，他始终摇摆不定，忽而宣讲它与天主圣言同永恒，忽而又说它由圣言的神性所造。因为他否认基督内有两个性体，一个神性的，一个人性的，一个来自圣父，一个来自其母亲，他认为圣言的本性被分割，仿佛一部分留在天主内，另一部分变成了血肉；因此，真理说两个性体中有一个基督，他却违反真理，断言基督的一个神性变成了两个性体。这就是\Person{阿波利纳里斯}的教义。\par

[35.] \Person{聂斯多略}的病症属于相反的一类，他假装在基督内承认两个不同的性体，却突然引入两个位格，并以闻所未闻的邪恶主张两个天主子、两个基督——一个神，另一个是人，一个由圣父所生，另一个由其母亲所生。因此，他坚持圣玛利亚不应称为\OriginalTerm{天主之母}{\GreekText{Θεοτόκος}}，而应称为\OriginalTerm{基督之母}{\GreekText{Χριστοτόκος}}，因为她所诞生的不是作为天主的基督，而是作为人的基督。但若有人以为在他的著作中他谈到一个基督，并宣讲一个基督的位格，切莫轻信。因为要么是狡诈的计谋，以便通过善更容易说服恶，正如宗徒所说：\Quote{\BibleQuote{那原是好的，竟成了我的死亡}} \BibleRef{罗 7:13}——要么，我再说，他在其著作的某些地方狡诈地装作相信一个基督和一个基督的位格，要么他说在贞女生育之后，两个位格联合成为一个基督，尽管在她受孕或分娩时以及此后短时期内，有两个基督；因此，固然基督起初生为一个普通的人，此外什么也不是，尚未与天主圣言在位格上合一，但后来圣言降下的位格降于他身上；尽管现在被接升的位格存在于天主的荣耀中，但曾经有一段时间他与所有其他人似乎并无差别。\par

% structure: p0072 source=h2 target=section reason=relative-heading-rank; base=h2

\section{第十三章}

\Emph{关于天主圣三与降生成人的天主教教义阐释}\par

[36.] 因此，这些狂犬\Person{聂斯多略}、\Person{阿波利纳里斯}和\Person{斐提努斯}以这些方式吠叫反对天主教信德：\Person{斐提努斯}否认天主圣三；\Person{阿波利纳里斯}教导圣言本性是可变的，并拒绝承认基督内有两个性体，此外还否认基督有灵魂，或至少否认他有理性灵魂，并断言天主的圣言取代了理性灵魂的位置；\Person{聂斯多略}则肯定一直有两个基督，或至少曾经有两个基督。但天主教会，对天主和对我们的救主都持有正确的信仰，无论是在天主圣三的奥迹中，还是在基督降生成人的奥迹中，都没有犯亵渎之罪。因为她在天主圣三的圆满中敬拜同一个神性，并在同一个至高尊严中敬拜圣三的平等，她承认一个耶稣基督，而非两个；同一个既是天主又是人，两者同样真实。她确实相信他内有一个位格，但两个性体；两个性体，但一个位格：两个性体，因为天主圣言并非可变以至于转化为血肉；一个位格，以免因承认两个儿子而显得她所敬拜的并非圣三，而是四位。\par

[37.] 但一再更清楚、更明确地阐述同一教义，将是合宜的。\par

在天主内，是一个性体，但三个位格；在基督内，是两个性体，但一个位格。在天主圣三内，位格与位格不同，而非性体与性体不同（位格有别，性体无异）；在救主内，性体与性体不同，而非位格与位格不同（性体有别，位格无异）。如何在天主圣三内，位格与位格不同，而非性体与性体不同？因为圣父是一个位格，圣子是另一个位格，圣神是又一个位格；然而，并非本性不同，而是同一本性。如何在救主内，性体与性体不同，而非位格与位格不同（两个不同的性体，而非两个不同的位格）？因为天主性是一个性体，人性是另一个性体。然而，天主性与人性并非两个不同的位格，而是同一个基督，同一个天主子，同一个基督和天主子的一位格，犹如在人内，肉体是一回事，灵魂是另一回事，但灵肉同属一人。在伯多禄和保禄内，灵魂是一回事，肉体是另一回事；然而，并非有两个伯多禄——一个是灵魂，另一个是肉体；或两个保禄——一个是灵魂，另一个是肉体；而是同一个伯多禄，同一个保禄，各自由两个不同的性体——灵魂与肉体——组成。同样，在同一个基督内有两个性体：一个是天主性，另一个是人性；一位（\OriginalTerm{出自}{ex}）天主圣父，另一位（\OriginalTerm{出自}{ex}）童贞圣母；一位与圣父同永恒、同等，另一位是受造的，次于圣父；一位与圣父同体，另一位与圣母同体；但在这两个性体中是同一个基督。因此，并非有一个基督是天主，另一个是真人；并非有一个是未受造的，另一个是受造的；并非有一个是不受苦难的，另一个是能受苦的；并非有一个与圣父同等，另一个次于圣父；并非有一个（\OriginalTerm{出自}{ex}）圣父，另一个（\OriginalTerm{出自}{ex}）圣母；而是同一个基督，天主而人，同样未受造和受造，同样不变且不受苦难，同样经历变化与苦难，同样与圣父同等和次于圣父，同样在时间之前由圣父所生（\Quote{在世界之前}），同样在时间之内由圣母所生（\Quote{在世界之内}），完全的天主，完全的人。在天主性内是至高的神性，在人内是完全的人性。我说完全的人性，因为它既有灵魂又有肉体；肉体是真正的肉体，是我们的肉体，是他母亲的肉体；灵魂是理性的，具有理智和理性。因此，在基督内有圣言、灵魂和肉体；但整体是一个基督，一个天主子，一个我们的救主和救赎者：唯一的一位，并非因天主性与人性某种败坏的混合，而是因位格的完整而独特的合一。因为结合并未将一性体转变或改变为另一性体（这是亚略派异端的典型错误），而是将两者如此结合为一，以致在基督内永远保持同一唯一位格的单一性，同时每一性体的特有属性也永远存留；由此得出，天主（即天主性）从未开始成为形体，形体也从未停止成为形体。这可在人性中说明：不仅在今世，而且在来世，每个人都将由灵魂和肉体构成；他的肉体永不会变成灵魂，灵魂也永不会变成肉体；但每个人将永远生活，两个性体之间的区别将在每个人内永远持续。同样，在基督内，每一性体将永远保持其特有属性，而无损于位格的合一。\par

% structure: p0077 source=h2 target=section reason=relative-heading-rank; base=h2

\section{第十四章}

\Emph{耶稣基督确是真人，非仅貌似}\par

[38.] 但当我们使用\Quote{位格}一词，并说天主藉一位格而成为人时，有理由担心我们的意思被理解为：天主圣言只是模仿地取了我们的人性，并履行人的行动，并非实在的人，而仅是貌似的人，犹如在戏台上，一个人在片刻间扮演几个角色，而他自己并非其中任何一个。因为当一个人承担扮演另一角色时，他履行该角色所应做的职分或行为，但他本人并非那个角色。因此，借用世俗生活中一个深受摩尼教徒喜爱的例子：当一名悲剧演员扮演司祭或君王时，他并非真正的司祭或君王。因为戏一结束，他所扮演的人物或角色便不复存在。但愿我们远离这等邪恶的嘲弄！让那迷惑人的宣讲者——摩尼教徒——的愚昧去吧，他们声称天主子、天主，并非真实的人，而是以虚假的言谈和生活方式假冒了人的角色。\par

[39.] 但公教信仰教导说，天主圣言降生成人是如此地真实：祂取了我们的本性，不是虚妄地、外表地，而是真实地、实际地，并行了人的行动，不是如同模仿别人的行动，而是作为祂自己的行动，并真实地就是祂所承担的那个人。正如我们自己，当我们说话、推理、生活、自立时，不是模仿人，而是人一样。例如，伯多禄和若望是人，不是通过模仿，而是真实地是人。保禄不是模仿宗徒，或假装自己是保禄，而是宗徒，就是保禄。同样，天主圣言在祂的屈尊中，通过取肉体并拥有肉体，借着肉体说话、行动和受苦，却无损于祂自己的天主性体，总而言之就是这一点：——祂不是模仿或假装自己是完人，而是显示自己真实地是真人。因此，如同灵魂与肉体结合，但没有变成肉体，不是模仿人，而是人，不是假装的人，而是真实的人；同样，天主圣言自己毫无变化，在与人结合时成为人，不是由于混淆，不是由于模仿，而是由于实际的存在和自立。那么，一次而永远地摒弃那种认为祂的位格是一个假装的角色的观念吧！在那里，总是存在的东西是一回事，假扮的东西是另一回事，扮演的人从来不是他所扮演的角色。天主不许我们相信天主圣言以这种虚幻的方式取了人的位格。我们宁可承认，祂自己的不变性体保持不变，同时取了完人的本性，祂自己实际上就是肉体，祂自己实际上就是人，祂自己实际上就是位格人；不是假装地，而是真实地；不是模仿地，而是本体地；最后，不是演出结束后就停止存在，而是实质地、永恒地存在和持续。\par

% structure: p0081 source=h2 target=section reason=relative-heading-rank; base=h2

\section{第十五章}

\Emph{天主性体与人性体的结合发生在童贞玛利亚的受孕中。称谓\Quote{天主之母}。}\par

[40.] 因此，基督位格的这种合一，不是在祂由童贞女出生之后才实现的，而是在她母胎中就已完成并完善的。因为我们必须特别注意宣认基督不仅是一位，而且永远是一位。因为如果你现在宣认祂是一位，却曾主张祂不是一位，而是两位，那是不可容忍的亵渎；即认为祂从受洗起是一位，但出生时是两位。除非我们承认人性与天主性体（借着位格的合一）相结合，不是从升天、复活或受洗开始，而是在祂母亲内，甚至在母胎中，甚至在童贞女受孕的那一刻，我们就绝对无法避免这种可怕的渎神。由于这种位格的合一，那些属于天主的属性归于人，那些属于肉身的属性归于天主，不加区别地、混同地。因此，在神圣引导下，一方面写道：人子从天降下（\BibleRef{若 3:13}），另一方面，光荣的主被钉在十字架上（\BibleRef{格前 2:8}）。由此，由于主的肉体被造成，主的肉体被创造，天主圣言本身也被说成被造成，全知的天主智慧也被说成被创造，正如先知性地描述祂的手和脚被刺透一样。从这种位格的合一，由于同样的奥迹，既然圣言的肉体由无玷的母亲所生，那么天主圣言本身由童贞女所生，这是最公教地相信的，最不敬神地否认的；既然如此，天主不许任何人试图剥夺圣玛利亚在神圣恩宠和特殊荣耀上的特权。因为由于那作为我们主和天主、同时也是她儿子的那一位的独特恩赐，她应该被最真实、最蒙福地宣认为——天主之母\Quote{\OriginalTerm{Theotocos}{\GreekText{Θεοτόκος}}}，但并不是像某种不敬神的异端所想象的那样，他们认为她被称为天主之母，仅仅因为她生了一个后来成为天主的男人，就如同我们称一个女人为司铎之母或主教之母，意思不是说她生了一个已经是司铎或主教的人，而是生了一个后来成为司铎或主教的人。我说，圣玛利亚并非这样成为\Quote{\OriginalTerm{Theotocos}{\GreekText{Θεοτόκος}}}，天主之母，而是如前所述，因为在她的圣胎中实现了那至圣的奥迹，由于位格独特而唯一的合一，正如圣言在肉体内是肉体，同样，人在天主内是天主。\par

% structure: p0084 source=h2 target=section reason=relative-heading-rank; base=h2

\section{第十六章}

\Emph{重述关于公教信仰和各种异端所说的话，第十一至十五章。}\par

[41.] 现在，为了重温我们先前关于上述异端或公教信仰的简要论述，让我们更简要、更简洁地再走一遍，以便重复后能更透彻地理解，并牢记不忘。\par

那么，福提诺该受诅咒，他不接受完整的天主圣三，却断言基督只是一个人。\par

阿波利纳里该受诅咒，他声称基督的天主性因转化而受损，并剥夺了祂完满人性的属性。\par

聂斯多略该受诅咒，他否认天主由童贞女所生，主张两个基督，并拒绝相信天主圣三，而引入四体。\par

但公教会是有福的，她在完整的天主圣三中敬拜一个天主，同时在天主性的合一中朝拜同等同荣的天主圣三，如此，既不会因本体的单一而混淆三位各自的特性，也不会因三位在圣三中的区别而分开天主性的合一。\par

我说，教会是有福的，她相信在基督内有两个真实而完善的性体，但只有一个位格，如此，既不会因性体的区别而分开位格的合一，也不会因位格的合一而混淆性体的区别。\par

有福的，我说，是那教会，她理解天主成为人，不是藉着本性的变换，而是因着位格的缘故，但这并非虚假易逝的位格，而是实在而永久的位格。\par

有福的，我说，是那教会，她宣称这\OriginalTerm{位格的合一}{unity of Person}是如此真实而有效，以致因着它，在一个\OriginalTerm{奇妙而不可言喻的奥迹}{marvellous and ineffable mystery}中，她把属神的属性归于人，并把属人的属性归于天主；因着它，一方面，她不否认那作为天主的\Person{人}从天降下，另一方面，她相信那作为人的\Person{天主}被造、受苦并被钉在十字架上；因着它，最后，她宣认\Person{人}是天主子，而\Person{天主}是童贞女之子。\par

那么，有福而可敬的，有福而至圣的，并且全然堪与天使天军的那天上赞颂相比的，是那用\OriginalTerm{三重圣颂}{threefold ascription of holiness}将荣耀归于独一主\Person{天主}的宣认。为此缘故，她还强有力地强调\OriginalTerm{基督位格的合一}{oneness of the Person of Christ}，以免她超越\OriginalTerm{天主圣三的奥迹}{mystery of the Trinity}（即实际上造成一个\OriginalTerm{四位一体}{Quaternity}）。\par

以上是题外话。他日若蒙主恩，我们将更详尽地处理并阐明这一主题。现在让我们回到正题。\par

% structure: p0096 source=h2 target=section reason=relative-heading-rank; base=h2

\section{第十七章}

\Emph{奥利振的错误是教会的一大试炼。}\par

[42.] 我们前面说过，在天主的教会中，教师的错误就是百姓的\OriginalTerm{试炼}{trial}，试炼的大小与犯错教师的学识成正比。我们首先用圣经的权威证明了这一点，然后举了教会历史中的实例，那些人一度被认为信仰纯正，最终要么落入某个已有的教派，要么创立自己的\OriginalTerm{异端}{heresy}。这确是一件重要的事，有益于学习，有必要记住，并且需要反复举例说明和强调，为使所有真正的天主教徒明白，他们应当与教会一起接受教师，而不是随从教师离弃教会的信仰。\par

[43.] 我相信，在可能列举的这类试炼的许多例子中，没有一个能与\Person{奥利振}的例子相比。在他身上有许多卓越、独特、令人钦佩之处，以致任何人理所当然地会认为他的一切断言都值得默然相信。\newline{} 若论言谈举止的权威，他勤勉非凡，谦逊、耐心、坚忍；若论出身或学问，还有什么比出身于因殉道而显赫的家族更高贵呢？后来，他为基督的缘故不仅失去了父亲，而且失去了所有财产，在神圣的贫困的窘迫中，他达到了如此高的标准，以至于据说他曾多次作为\OriginalTerm{宣信者}{Confessor}受苦。\newline{} 与他有关的还不止这些，后来这一切都成了试炼的缘由。他的天赋如此有力、如此深邃、如此敏锐、如此优雅，几乎没有人能远远超越他。他的学识和一般学问的光辉如此灿烂，以至于神圣哲学中几乎没有几点、世俗学问中几乎没有哪一点是他没有彻底掌握的。当希腊文屈从于他的勤奋时，他又精通了希伯来文。\newline{} 关于他的口才，我还要说什么呢？其文风如此迷人、柔和、甜美，以致从他口中流出的仿佛是蜂蜜而非言语！有什么课题，无论多么困难，他没有以其推理的力量使之清晰明了呢？有什么事业，无论多么艰难，他没有使之显得极其容易呢？\newline{} 但也许他的论断仅仅基于巧妙编织的论证？恰恰相反，从未有教师像他那样引用如此多的圣经证据。那么，我想他著作不多？没人比他更多，所以，若我没有记错，他的著作不仅无法全部读完，甚至无法全部找到；因为为了不给他的知识机会留下任何欠缺，他还额外享有长寿的恩惠。\newline{} 但也许他在门徒方面并不特别幸运？谁比他更幸运？从他的学派中走出了无数博士、司铎、宣信者、殉道者。\newline{} 那么，谁能表达他如何被众人敬仰？他的名望有多大？他的影响有多广？有谁的宗教情操在普通水准之上却不从地极赶来看他？有哪个基督徒不几乎像尊敬先知一样尊敬他？有哪个哲学家不把他当作大师？\newline{} 不仅普通人，连宫廷对他的崇敬有多大，那些讲述他如何被\Person{亚历山大皇帝}的母亲派人召去的历史可以证明，她因他所激发起的对她所热爱的天上智慧的爱而感动。他的信件也为此作证，他以基督教导师的权威写信给\Person{腓力皇帝}——第一位信奉基督教的罗马王子。\newline{} 至于他令人难以置信的学问，如果有人不愿接受我们基督徒的证言，那么至少让他接受哲学家中异教徒的证言。因为那不敬虔的\Person{波斐利}说，当他还是个男孩时，因仰慕他的名声而去了\Place{亚历山大港}，在那里看到了他，那时他已年老，但显然是一位造诣如此之深的人，以致他已达到了普遍知识的顶峰。\par

[44.] 纵使我只略述少许，时间也不够我讲述此人的\OriginalTerm{卓越}{excellencies}，而这一切不仅有助于宗教的光荣，也增加了这\OriginalTerm{试炼}{trial}的严重性。因为在世上，有谁会轻易舍弃一位拥有如此伟大天赋、如此学识、如此影响的人，而不更愿意采用那句名言：宁可跟随\Person{奥利振}犯错，也不与他人正确？\par

我还要说什么呢？结果是，许多人被引离了信德的完整，不是因为这位伟人、这位伟大导师、这位伟大先知的任何人性卓越，而是——正如事件所表明的——因为他所证明的过于危险的考验。于是，这位奥力振，尽管如此伟大，却妄用天主的恩宠，轻率地随从自己天才的趋向，过分自信，轻视基督信仰古老的纯朴，自以为知道得比世上所有人都多，藐视教会的传统和古人的决断，以新奇的方式解释了一些圣经段落，因而为自己招来了给予天主教会的那项警告，它适用于他如同适用于他人：\BibleQuote{如果有先知在你们中间兴起，}……\BibleQuote{你们不可听从那先知的话，}……\BibleQuote{因为上主你们的天主试探你们，要试验你们是否爱祂。}\BibleRef{申 13:1} 的确，如此突然地诱惑那全心崇拜他、因仰慕他的天才、学问、口才、生活方式和影响力而依恋他的教会——她毫无畏惧、毫不猜疑——从古老的宗教慢慢、逐渐地引向新的亵渎，这不仅是一个考验，而且是一个大考验。\par

[45.] 但有人会说，奥力振的书籍被篡改了。我不否认；不仅如此，我欣然认同。因为这种情况已通过口传和文字流传下来，不仅由天主教徒，也由异端者。但问题在于，尽管他本人不这样，但以他名义出版的书籍却成了一个大的考验，这些书籍充满许多有害的亵渎之言，被阅读并喜爱，不是作为别人的作品，而是被认为是他的作品，因此，即使奥力振的原意没有错误，奥力振的权威似乎成了引导人们接受错误的有效原因。\par

% structure: p0103 source=h2 target=section reason=relative-heading-rank; base=h2

\section{第十八章}

\Emph{戴尔都良：教会的大考验}\par

[46.] 戴尔都良的情况也是如此。因为正如奥力振在希腊人中独占鳌头，戴尔都良在拉丁人中也是如此。有谁比他更有学问，有谁比他更精通神知或人知？他以非凡的心智能力通晓一切哲学，知道所有哲学学派以及学派的创始人和拥护者，熟悉他们的一切规则和遵守，以及他们各种历史和学问。他的天才不是如此无与伦比的强大和猛烈，以致几乎没有任何他设想要克服的障碍，不是被他以敏锐刺穿，就是以重量压垮？至于他的文体，谁能充分赞扬？其论证如此有力，以致即使在未能说服时，也迫使同意。几乎每个词都是一段警句；每句话都是一场胜利。马尔基翁、阿培勒斯、普拉克塞亚斯、赫尔谟根斯、犹太人、异教徒、诺斯替派以及其余的人都知道这一点，他们的亵渎被他用众多巨著的力量像雷电一样击倒。然而，尽管我提到了这一切，这位戴尔都良，我说，对公教教义，即对普世而古老的信仰坚持得太过脆弱，远更雄辩而非忠诚，改变了他的信仰，并证实了真福精修者希拉利关于他所写的话，即 \Quote{他随后的错误损害了他认可著作的权威。} 他在教会中也是一个大的考验。但我不愿对戴尔都良多说什么。我只补充一点：违背梅瑟的禁令，他断言在教会中兴起的蒙丹的新癫狂，以及疯女人们梦见的关于新学说的疯狂梦魇，是真实的预言，因而他和他的著作都配得上这样的话：\BibleQuote{如果有先知在你们中间兴起，}……\BibleQuote{你们不可听从那先知的话。} 因为为何？\BibleQuote{因为上主你们的天主试探你们，要试验你们是否爱祂。}\par

% structure: p0106 source=h2 target=section reason=relative-heading-rank; base=h2

\section{第十九章}

\Emph{我们应从这些例子学到什么}\par

[47.] 那么，我们应当留意教会历史上这些如此众多且重大的例子，以及其他类似的例子，并清楚理解，按照\BibleBook{}{申命纪}中制定的规则，如果教会的某位导师任何时候在信德上犯了错误，天主眷顾允许它发生，为要试探我们，是否我们全心、全意爱天主。\par

% structure: p0109 source=h2 target=section reason=relative-heading-rank; base=h2

\section{第二十章}

\Emph{真正天主教徒的特征}\par

[48.] 既然如此，那真正而纯正的天主教徒是爱天主的真理、爱教会、爱基督的身体、看重神圣的宗教和公教信德胜过一切，胜过任何人的权威、尊重、天才、口才、哲学；他轻视这一切，并持续坚定、稳固于信德，决心只相信他确信天主教会从古至今普世所持守的；但凡他发现有人偷偷引入的、有别于或违反所有圣人所传的任何新奇而前所未闻的教义，他会明白这不属于宗教，而是作为考验被允许，尤其受蒙真福宗徒保禄的话教诲，他在\BibleBook{格林多}{前书}中写道：\BibleQuote{要有异端，好使那经过考验的人在你们当中显明出来。}\BibleRef{格前 2:9} 好像他说，这就是为何异端的创始者没有被天主立即拔除的原因，即为了使那经过考验的人显明出来；也就是说，使每个人显明出来，看他对公教信德的爱多么坚定、忠信而稳固。\par

事实上，每当出现一种新奇事物，立刻就能辨出麦粒的重量与糠秕的轻浮之间的差别。那时，那在场上没有重量支撑的东西就轻易被吹走了。因为有些人立即完全飞离；另一些人只是被摇动，害怕灭亡，受伤，半死不活，羞于返回。事实上，他们吞下了不少毒药——不足以致死，却又多于能排除的；它既不致死，也不容许活着。悲惨的境况！他们被何等汹涌波涛般的忧思所抛掷！一时，错误被启动，他们随风飘荡；一时，又像回流的海浪折返，被冲回；一时，轻率地赞同那看似不确定之事；一时，又不理性地惧怕那确定之事，不知去向何处、折返何处、寻求什么、躲避什么、保持什么、抛弃什么。\par

确实，这种犹豫不决、悲惨动摇的心灵的苦难，如果他们明智的话，是天主怜悯为他们预备的一剂良药。为此，在天主教信仰最安全的港湾之外，他们被各种动荡的思绪抛掷、击打、几乎丧命，为的是他们能收起那自满的帆——那帆他们曾不智地向着新奇的风暴张开——然后投奔归于他们温和而善良的母亲那最安全的港湾，并停留在那里，开始吐出他们曾吞下的那些苦涩而浑浊的错误洪流，以便此后能饮用新鲜活水的水流。让他们好好忘记他们曾学得不好的东西，并让他们领受教会全部教义中他们能理解的那么多：他们不能理解的，让他们相信。\par

% structure: p0114 source=h2 target=section reason=relative-heading-rank; base=h2

\section{第二十一章}

\Emph{圣保禄言论释义——\BibleRef{弟前 6:20}}\par

既然如此，当我思量这些事，并在心中反复思索时，我不禁对某些人的疯狂、他们盲目理解的邪恶、以及对错误的贪恋感到万分惊奇，因为他们不满足于那一次交付并从古时领受的信仰准则，而是日日寻求一个接一个的新奇，不断渴望在宗教中增添、改变、删减，仿佛\Quote{一旦启示了的，就当满足}这教训不是属天的准则，而是属地的——一个只有通过不断的修正，甚至不断的指摘才能遵守的准则。然而，神圣的训谕大声呼喊：\BibleQuote{不可挪移你祖先所立的界碑}\BibleRef{箴 22:28}，\BibleQuote{不要与审判官争讼}\BibleRef{德 8:14}，\BibleQuote{谁拆毁篱笆，蛇必咬他}\BibleRef{训 10:8}，以及宗徒的那句话，它如同属灵的利剑，曾斩断一切异端的所有邪恶新奇，也必将永远斩断：\BibleQuote{弟茂德啊，要保管所受的寄托，躲避亵渎的新奇言论和伪知识相反的主张，有些人自负有这种知识，就在信仰上走差了。}\BibleRef{弟前 6:20}\par

听了这样的话，还有谁如此铁石心肠、如此顽梗、如此如钻石般固执，以至于不被如此巨大的打击所屈服，不被如此沉重的重量所压碎，不被如此重的击打所震碎，不被如此可怖的神圣言辞的雷霆所消灭呢？\Quote{躲避亵渎的新奇言论}，他说。他没有说\Quote{躲避古老}。但他清楚地通过相反的原则指明了应当遵循什么。因为如果新奇应当被躲避，那么古老就应当被持守；如果新奇是亵渎的，那么古老就是神圣的。他接着说：\Quote{和伪知识相反的主张。}\Quote{伪知识}的确，指异端者的学说，那里无知伪装成知识，迷雾伪装成阳光，黑暗伪装成光明。\BibleQuote{有些人自负有这种知识，就在信仰上走差了。}自负什么？不过是某些（我不知道是些什么）新而从未听过的学说。因为你可能听到这些所谓博士中有人说：\Quote{来，愚蠢的可怜虫，你们被称为天主教徒，来学习真正的信仰，这信仰除我们之外无人知晓，它隐藏了许多世代，但最近被揭示和显明。但要偷偷地学，秘密地学，因为你们会喜欢它。此外，一旦你们学会了，要秘密地教导，免得世界听见，免得教会知道。因为只有少数人蒙恩接受如此伟大奥秘的秘密。}这不正是那妓女在\BibleBook{}{箴言}中向行路者呼喊的话吗？\BibleQuote{谁是无知的，请转身到这里来。}至于那些缺乏领悟的人，她劝诫他们说：\BibleQuote{偷来的水是甜的，暗吃的饼是美味的。}接下来呢？\BibleQuote{他却不知道阴魂在她家中沉沦。}\BibleRef{箴 9:16}谁是那些\Quote{阴魂}？让宗徒来解释：\Quote{那些在信仰上走差了的人。}\par

% structure: p0118 source=h2 target=section reason=relative-heading-rank; base=h2

\section{第二十二章}

\Emph{对弟茂德前书第六章二十节更详细的释义\BibleRef{弟前 6:20}}\par

但是，全面阐释宗徒的那段话是值得的：\BibleQuote{弟茂德啊，要保管所受的寄托，躲避亵渎的新奇言论。}\BibleRef{弟前 6:20}\par

\Quote{啊！}这感叹既有预知，亦有爱德。因为他预见到将要发生的错误而提前哀恸。今日的弟茂德是谁？或者是普世教会，或者特指整个主教团，他们自己应当拥有完备的宗教知识，并将之传授与他人。所谓\Quote{保管所托付的真理}是什么？\Quote{保管它}，因有盗贼，因有仇敌，免得在人睡觉时，他们在人子撒在田里的好麦子上撒下莠子。\Quote{保管所托付的真理}。什么是\Quote{所托付的真理}？那是交托给你的，而非你自己发明的；是关乎学识，而非才智；是公共传承，而非私人采纳；是带给你的，而非你提出的；你在此不是作者，而是保管者；不是师傅，而是门徒；不是领袖，而是跟随者。\Quote{保管所托付的真理}。要保全天主教信德的才能，完整无缺，纯净无杂。那交托给你的，让它存留于你，由你传递。你接受了金子，就要还以金子。不可用一物代替另一物。不可无耻地用铅或铜代替金子。要交出真金，而非赝品。\par

啊！弟茂德！啊！司铎！啊！诠释者！啊！导师！若天赐的恩宠使你有才智、技巧、学识，就当成为属灵帐幕中的\Person{贝匝肋耳}，雕刻神圣教义之珍贵宝石，精确镶嵌，巧妙装饰，增添光彩、恩宠、美丽。让那先前仅被相信却未完全领悟的真理，经由你的诠释而得以清晰理解。让后代因着你的阐述而欣然接受古时所敬拜却未理解的道理。然而，你所教导的仍当是你所学到的同一真理，以致你虽以新方式宣讲，所言却非新事。\par

% structure: p0123 source=h2 target=section reason=relative-heading-rank; base=h2

\section{第二十三章}

\Emph{论宗教知识的进展}\par

[54.] 但或许有人会说：那么，在基督的教会中就没有进步了吗？当然有，而且应有尽有。因为，有什么样的人会如此嫉妒他人、如此憎恨天主，竟想禁止进步呢？但条件是：这必须是真正的进步，而非信仰的改变。因为进步要求主题本身得以扩展，而改变则是使之转化为别的东西。因此，无论是个人的还是全教会的理智、知识、智慧，都应当在时代和世纪的过程中增长并取得巨大而有力的进步；然而，这进步只能在其本类之内，也就是说，在同一教义、同一意义、同一含义中。\par

[55.] 宗教在灵魂中的成长必须类比身体的成长：身体虽随岁月发展并达到其完满尺寸，但仍保持同一。青春之花与成熟之年有巨大差别；然而曾经年轻的人，当他们变老时，仍然是同一个人，尽管个人的身材和外表改变了，但其本性是同一个，其位格是同一个。婴儿的肢体很小，年轻人的很大，但婴儿和年轻人是同一个人。成年人与儿童拥有相同数量的关节；若有任何成员是成熟年龄才产生的，它们早已以胚胎形式存在，因此，老年时并无新物产生，而只是在孩童时已潜伏的东西。因此，这无疑是真实合法的进步法则，是固定而最美的成长秩序：成熟年龄在人身上发展那些造物主的智慧早已在婴儿中预先塑造的部分和形式。反之，如果人的形状变成另一种类，或者肢体数量增加或减少，那么整个身体就会变成残废或怪物，或至少受损衰弱。\par

[56.] 同样，基督宗教的教义也应当遵循同样的进步法则：被年岁巩固，被时间扩展，被年龄精炼，然而始终保存完整无缺、纯净无杂，在其所有部分的度量中完备而完美，并且可以说，在其所有适当的肢体和意义上，不容许任何改变、任何特性的损失、任何界限的变动。\par

[57.] 例如：我们的祖先在古时曾在教会的田里撒下麦子。若我们这些后代，竟收获莠子的虚假错误，而非谷物的真正真理，那是极不相称、极不公义的。相反，应当如此：最初与最后之间不应有差异。从作为麦子撒下的教义，我们应当在增长中收获同类教义——也是麦子；这样，当随着时间的推移，最初的种子得以发展并因栽培而繁荣时，植物的特性不得改变。外形、形式、外表的变化可以出现，但每种本性的本质必须保持不变。但愿那些天主教诠释的玫瑰花圃不要变成荆棘和蒺藜。但愿在那属灵的乐园里，肉桂和香脂的植物中，不要突然长出毒麦和乌头。\par

因此，凡是由教父的忠信在这天主管家的耕作中所撒下的，都应当由子孙的勤勉加以培育和照料，都应当繁荣成熟，都应当前进并走向完美。因为，那些天上哲学的古老教义，应当随着时间推移而得到照料、润饰和打磨；但不应被改变、被截肢、被肢解。它们可以接受证明、说明、确定；但必须保持其完整性、纯正性、特性。\par

[58.] 因为一旦允许了这种不敬的欺骗自由，我害怕说出宗教将面临多么巨大的危险，以致完全毁灭和消失。因为如果公教真理的任何一部分被放弃，那么另一部分、另一部分、又一部分将理所当然地随之被放弃，而当各个部分被拒绝后，最终除了全盘拒绝还会有什么结果？另一方面，如果新事物开始与旧事物相混，外来的与本土的相混，亵渎的与神圣的相混，那么这种习俗必然普遍蔓延，直到最后教会将没有任何未经篡改、没有任何纯正、没有任何健全、没有任何纯洁的东西；而是从前有贞洁无瑕真理的圣所，从此却成了不敬和卑劣错误的妓院。愿天主的仁慈将这邪恶从祂仆人的心中除去；愿它成为不敬者的疯狂。\par

[59.] 然而基督的教会，她是委托给她保管的教义的谨慎而警醒的守护者，从未在其中改变任何东西，从未减少，从未添加，未剪除必要的，未增加多余的，未失去自己的，未侵占他人的；而是在忠实而明智地处理古老教义的同时，小心地保持这一目标——如果古代留下的任何东西是未成形和初步的，就去塑造和打磨它；如果任何东西已经成形和发展，就去巩固和加强它；如果任何东西已经批准和界定，就去保持和守护它。最后，大公会议在其法令中还曾追求过什么别的目标，除了确保以前单纯相信的将来要明智地相信，以前冷淡宣讲的将来要热诚地宣讲，以前疏忽实行的将来要加倍谨慎地实行？我说，这就是公教会，被异端的新奇事物所激发，通过其大公会议的法令所完成的——这就是，而且仅此而已——她从此将只以传统从古人那里领受的东西以书面形式传给后代，用少量言辞包含大量内容，并且常常为了更好地理解，用新名称的特征来指定古老的信仰信条。\par

% structure: p0132 source=h2 target=section reason=relative-heading-rank; base=h2

\section{第二十四章}

\Emph{延续对\BibleRef{弟前6:20}的阐释}\par

[60.] 但让我们回到宗徒。他说：\Quote{啊，弟茂德，} \Quote{要保管好寄托，避开亵渎言语的新奇。} \Quote{要避开它们，如同避开毒蛇、蝎子、毒蜥，免得它们不仅用触摸，甚至用眼睛和气息击打你。} 什么是\Quote{避开}？甚至不与这种人\BibleRef{格前5:11}一起吃饭。什么是\Quote{避开}？圣若望说：\Quote{若有人来到你们那里，不带这教义，} 什么教义？除了那公教与普世的教义，它通过真理的未腐传统在几个时代的延续中始终如一，并将永远如此—— \Quote{不要接他到家里，也不要向他问安；因为向他问安的，就在他的恶行上有份。} \BibleRef{若二1:10}\par

[61.] \Quote{亵渎言语的新奇。} 这些是什么言语？那些毫无神圣、毫无宗教的言语，完全远离教会最内在的圣所——即天主的圣殿。亵渎言语的新奇，即教义、主题、意见的新奇，与古老和古时的信仰相悖。如果这些被接受，那么必然的结果是：有福教父们的信仰要么全部，至少大部分被违反；必然的结果是：所有时代的信徒、所有圣人、贞洁者、节欲者、贞女、所有圣职人员、执事和司铎、成千上万的宣信者、如此庞大的殉道者军队、如此众多的城市和民族、如此多的岛屿、省份、君王、部落、王国、列国，总之，几乎整个大地，通过公教信仰与元首基督结合，在如此漫长的时间里一直无知、错误、亵渎，不知道相信什么、宣认什么。\par

[62.] \Quote{避开亵渎言语的新奇，} 接受和追随这些从来不是公教徒的事，而是异端者的事。事实上，有什么异端不是以特定的名称、在特定的地点、特定的时间爆发的？有谁开创了异端，却没有首先脱离公教会普遍性与古老性的一致同意？这一点通过例子得到了最清楚的证明。因为在亵渎的\Person{白拉奇}之前，有谁曾将如此大的先天力量归于自由意志，以至于否认天主的恩宠在每个行动中帮助它向善的必要性？在他那怪异的门徒\Person{塞莱斯提乌}之前，有谁曾否认全人类都沾染了亚当罪的罪债？在亵渎的\Person{亚略}之前，有谁胆敢撕裂天主圣三的合一？在邪恶的\Person{撒伯流}之前，有谁如此鲁莽，去混淆天主圣三的一体？在残忍的\Person{诺瓦替安}之前，有谁将天主描绘为残忍，宁可恶人死亡而不愿他悔改得生？在\Person{西满术士}之前——他被宗徒的斥责击打，从他那里一切卑劣事物的古老源头通过隐秘的连续传承一直流到我们时代的\Person{普里西利安}——谁，我说，在这\Person{西满术士}之前，胆敢说创造者天主是邪恶的创造者，即我们的恶行、不敬、丑行的创造者？因为他断言祂亲手创造了这样一种人性，它凭自身运动并在其必然约束意志的冲动下，除了犯罪不能做任何事，不能愿意任何事，因为它被各种恶习的狂怒摇撼和点燃，被无法熄灭的贪欲拖入卑劣的极端。\par

[63.] 这类例证多不胜数，为求简洁，不一一列举；然而，这一切都显然昭示一条确立的法则：几乎在所有异端中，他们永远喜好亵渎性的新奇，藐视古人的决断，并因虚假知识的反对，使信德遭逢船难。另一方面，天主教徒的确切标志是：持守圣教父所托付的，谴责亵渎性的新奇，并用宗徒再三重复的话，咒诅任何传播与所领受的不同教导的人。\BibleRef{迦 2:9}\par

% structure: p0138 source=h2 target=section reason=relative-heading-rank; base=h2

\section{Chapter 25.}

\Emph{异端引用圣经，为更容易得逞于欺骗。}\par

[64.] 这里，或许有人会问：异端也引用圣经吗？他们确实引用，而且变本加厉；因为你会看到他们如飞疾走，穿过圣经的每一卷书——梅瑟的书、列王纪、圣咏集、书信、福音、先知书。不管在自己人中，还是在外人面前，私下或公开，言谈或写作，在宴饮聚会时，或在大街上，他们几乎从不提出任何自己的东西，而不企图以圣经的话来庇护。阅读\Person{撒莫萨塔的保禄}、\Person{普里西利安}、\Person{欧诺米}、\Person{约维尼安}以及其余那些害虫的作品，你会看到数不清的实例，几乎没有一页不布满来自新约或旧约的看似有理的引文。\par

[65.] 但他们越是隐秘地隐藏于神圣法律的庇护之下，就越应当被畏惧和警惕。因为他们知道，他们教义的恶臭，如果赤裸裸地散发出来，几乎不会为任何人所接受。因此，他们以此而洒上神圣语言的香水，致使一个本来会轻易鄙视人的错误的人，可能犹豫于谴责神圣的话语。事实上，他们所做的，正如乳母准备给孩子的苦药时所做的：她们在杯口四周涂上蜂蜜，使毫无戒心的孩子先尝到甜味，便不害怕苦味。这些人也是如此，他们用药物之名来掩饰毒草和毒汁，几乎无人，一看标签，就怀疑那是毒药。\par

[66.] 正因如此，救主呼喊说：\BibleQuote{你们要提防假先知，他们来到你们这里，外面披着羊皮，里面却是凶残的狼。} \BibleRef{玛 7:15} 所谓\Quote{羊皮}是什么意思？除了先知和宗徒们以羊的纯真所预先编织的话，作为那除去世界罪恶的无玷羔羊的羊毛，还能是什么？凶残的狼又是什么？除了异端不驯和狂野的曲解，他们不断侵扰教会的羊栈，凡力所能及，撕裂基督的羊群，还能是什么？但他们要更成功地用诡计偷袭没有戒备的羊群，保留狼的凶残，却脱下狼的外表，以神圣法律的话包裹自己，就像裹在羊毛里，以至于人一摸到羊毛的柔软，就不害怕狼的尖牙。但救主说什么？\Quote{凭他们的果实，你们能认出他们来。} 也就是说，当他们开始不仅引用这些神圣的话，而且开始解释它们，不单以此夸耀站他们一边，也开始诠释时，那时，那苦涩、那尖刻、那愤怒将被理解；那时，那新奇毒药的恶臭将被人察觉；那时，那些亵渎性的新奇将被揭露；那时，你会首先看到篱笆被破坏，继而是教父的界标被挪移，继而公教信德被攻击，继而教会教义被撕裂。\par

[67.] 这些人就是宗徒保禄在致格林多人后书中所责斥的，他说：\BibleQuote{因为这种人，是假宗徒，是诡诈的工人，假装成基督的宗徒。} \BibleRef{格后 11:12} 宗徒们从圣经中引出例子，这些人也同样；宗徒们引用圣咏集的权威，这些人也同样；宗徒们引用先知书的话，这些人还是同样。但当他们开始用不同意义来解释双方都同意引用的段落时，纯真者与诡谲者、真品与赝品、正直者与扭曲者就辨别出来了，总之，真宗徒与假宗徒就辨别出来了。\Quote{这并不奇怪，} 他说，\BibleQuote{因为撒殚自己也假装成光明的天使。所以，如果他的仆人假装成正义的仆人，也不算什么大事。} 因此，根据宗徒保禄的权威，每当假宗徒或假教师引用神圣法律的话，并通过错误解释来支撑自己的错误时，毫无疑问，他们在追随他们父亲的诡计；这诡计，他肯定绝不会想出，除非他知道，在他能用欺诈和隐秘引入错误的地方，没有比假装圣经的权威更容易实现其不敬目的的办法了。\par

% structure: p0144 source=h2 target=section reason=relative-heading-rank; base=h2

\section{Chapter 26.}

\Emph{异端引用圣经，是效法魔鬼的榜样。}\par

[68.] 但有人会说，我们有什么证据证明魔鬼惯于引用圣经？让他读福音，其中记载：\BibleQuote{那时，魔鬼就带他到耶路撒冷，叫他站在圣殿顶上，对他说：\Quote{你若是天主子，就跳下去，因为经上记载：\InnerQuote{他为你吩咐自己的天使，用手托着你，免得你的脚碰在石头上。}}} 那么，像我们这样微不足道的可怜人，该受那竟敢用经文攻击荣耀之主的人怎样的对待呢？他说：\BibleQuote{你若是天主子，就跳下去。} 为什么？他说：\BibleQuote{因为经上记载。} 我们应当特别注意这段经文并铭记在心，因受如此重要的福音权威的警告，我们可以确信无疑：当我们发现有人引用宗徒或先知的话反对公教信仰时，那是魔鬼通过他们的口说话。因为正如那时元首对元首说话，如今肢体对肢体说话：魔鬼的肢体对基督的肢体，妄信者对信徒，亵渎者对虔诚者，一句话，异端者对公教徒。\par

[69.] 他们说什么？\BibleQuote{你若是天主子，就跳下去；} 也就是说，如果你愿作天主子，并愿承受天国的产业，你就跳下去；也就是说，从这崇高教会（人们以为无非就是天主圣殿）的教导和传统中跳下去。如果有人问这些献计者中的一个异端者：你如何证明？你有什么根据说，我应当抛弃公教会那普遍而古老的信仰？他早已准备好答案：\BibleQuote{因为经上记载。} 并且他立刻提出上千个证据、上千个事例、上千个权威，取自法律、圣咏、宗徒、先知，通过这些按新而错误的原则解释的经文，那不幸的灵魂便可能从公教真理的高峰跌入异端的深渊。然后，伴随着应许，异端者惯于奇妙地欺骗粗心者。因为他们敢教导并承诺，在他们的教会中，也就是在他们的聚集中，有一种伟大而特殊且完全个人性的天主的恩宠，以至于凡属于他们一伙的，无需任何劳苦、无需任何努力、无需任何勤奋，甚至不求、不寻、不敲门，就获得天主如此的安排，以至于被天使之手托着，即受天使的保护，他们永远不可能用脚碰在石头上，即永远不会跌倒。\par

% structure: p0148 source=h2 target=section reason=relative-heading-rank; base=h2

\section{第二十七章}

\Emph{解释圣经时应遵守什么规则}\par

[70.] 但有人会说，如果魔鬼及其门徒——其中一些是假宗徒，一些是假先知和假教师，他们无一例外都是异端者——都引用圣经的言语、思想和应许，那么作为公教徒和慈母教会的子女该做什么？他们如何在神圣圣经中分辨真理与谎言？他们必须非常谨慎地遵循我们在本备忘录开头所提及的圣洁而博学的人所推荐的方法，也就是说，他们必须根据普世教会的传统并遵循公教教义的规则来诠释神圣正典，此外，在公教和普世教会中，他们必须遵循普遍性、古老性和一致性。如果任何时候，部分反对整体，新奇反对古老，一个或少数犯错者的异议反对所有或至少大多数公教徒的一致意见，那么他们应当首选整体的健康而非部分的腐败；在这个整体中，他们应当首选古老宗教的新奇亵渎；同样，在古老性本身中，对于一个人或极少数人的鲁莽，他们应当首先首选普世大公会议的总决定（如果有的话），如果没有，则其次应遵循许多伟大教父的一致信念。如果忠实、冷静和谨慎地遵守这一规则，我们就能轻而易举地察觉并消除异端者有害的谬误。\par

% structure: p0151 source=h2 target=section reason=relative-heading-rank; base=h2

\section{第二十八章}

\Emph{如何通过比较古代教父的一致意见来察觉并谴责异端的新奇之说}\par

[71.] 在此，我意识到，作为上述内容的必然结果，我应当举例说明如何通过比较古代教父的一致意见来察觉并谴责异端者的亵渎新奇之说。然而，在调查圣教父们的这种古老一致意见时，我们不应将精力花费在神圣法律的每一个次要问题上，而应只（至少特别）关注涉及信仰准则的问题。而且，这种处理异端的方法不应总是或每一种情况下都采用，而只适用于那些新近出现的异端，并且在它们初起之时，在它们还未有时间败坏古老信仰的准则之前，在它们试图随着毒素扩散和蔓延而腐蚀古代文献之前。但是对于那些已经广泛传播且历史悠久的老异端，决不可如此处理，因为随着时间流逝，它们早已有机会败坏真理。因此，对于更古老的分离或异端，我们应当要么（如有必要）仅凭圣经的权威来驳斥它们，要么至少回避它们，因为早在古代，它们就已经被天主教司铎的普世公会议定罪和谴责了。\par

[72.] 因此，每当任何有害错谬的败坏开始显露，并藉窃取某些圣经章节、以欺诈诡诈的方式加以解释来自我辩护时，就应立即收集古代教父在诠释圣经正典时的意见，以便使那兴起的异说——无论它是什么——既无可置疑地被揭露，又无可推诿地受谴责，并因而显出其亵渎性。但只有那些在公教信仰与共融中圣洁、明智且恒心地生活与教导，并堪当在基督信仰中去世或为基督欣然殉道的教父的意见，才可用于比较。然而，我们相信他们时须持此条件：只有那些所有教父或大多数教父在同一意义上明确、频繁、持续地支持并确认的意见——如同形成教父们的一致会议，众人接受、持守、传授同一道理——才被视为无可置疑、确定、稳固的。但凡某位教师所持的异于众人或与众人相反的意见，无论他多圣洁、多有学识，无论他是主教、宣信者还是殉道者，都应视为他个人的私见，应与共同、公开、普遍信念的权威分离，以免我们追随一人新创的错谬而危及永恒的救恩，如同异端者和分裂者亵渎的惯行，抛弃普遍信经的古老真理。\par

[73.] 恐怕有人会轻率地认为这些蒙福教父之神圣公教的一致意见可以轻视，宗徒在致格林多人前书中说：\Quote{天主在教会内所设立的：第一是宗徒}，\BibleRef{格前 12:27}，他自己便是其中之一；\Quote{第二是先知}，如我们在宗徒大事录中所读的阿加波，\BibleRef{宗 11:28}；\Quote{第三是教师}，即现今称为讲道者、讲解者的人，同一位宗徒有时也称他们为\Quote{先知}，因为藉着他们，先知书的奥秘向民众开启。因此，凡轻视这些由天主在教会中按各时期、各地域所设立的人，当他们为解释公教道理的某一点而在基督内意见一致时，他所轻视的便不是人，而是天主；为使无人偏离真理中的合一，同一位宗徒恳切告诫：\Quote{弟兄们，我因我们的主耶稣基督之名，求你们众人言谈一致，在你们中间不要有分裂，但要同心合意，完全相合。}\BibleRef{格前 1:10} 但若有人不同意他们一致的决定，愿他听从同一位宗徒的话：\Quote{天主不是混乱的天主，而是平安的天主；}\BibleRef{格前 14:33} 即，祂不是那脱离一致平安者的天主，而是那在一致平安中持守者的天主；\Quote{正如我在各处圣徒的教会中所教训的}，即公教会的教会；这些教会之所以是圣徒的教会，是因为她们在信仰的共融中恒心持守。\par

[74.] 并且，为防任何人藐视众人、傲慢地要求单听他自己、单信他自己，宗徒接着说：\Quote{难道天主的道理是从你们出来吗？或是单临到你们吗？}又为免此话被视为轻率，他继续：\Quote{若有人自以为是先知，或是属神的人，就该承认我所写给你们的是主的命令。}对此，除非一个人是先知或属神的人，即属灵之事上的导师，否则他应尽可能持守公正与合一，既不偏私自己的意见超过众人，也不背离整个团体的信仰。凡忽略这训令的，他将被忽略；\BibleRef{格前 14:33} 意思是，凡不学习他所不知道的，或藐视他所知道的，他必被忽略，即被视为不配与那些藉信德彼此联合、藉谦卑彼此平等的人同列于天主面前——我想象不出比这更可怕的灾祸。然而，这正是我们所见到的，按照宗徒的警告，临到了佩拉基乌斯派的尤利安身上：他或忽视了与同道基督徒的信仰认同，或妄自与之分离。\par

[75.] 但现在该提出我们应许过的例证，即何时、何处曾收集圣教父的判决，以便按照它们，藉由大公会议的谕令与权威，确立教会信仰的准则。为使此事更便于进行，且让这本《备忘手册》在此结束，以便接下来的部分能从新的起点开始。\par

\EditorialNote{《备忘手册》的第二卷已失传，仅存其结论，即随后的总结。}\par

% structure: p0159 source=h2 target=section reason=relative-heading-rank; base=h2

\section{第二十九章}

\Emph{总结。}\par

[76.] 既然如此，现在是在这第二卷《备忘手册》结束时，将本卷与前卷所述之内容加以总结的时候了。\par

我们此前曾说，公教徒历来并至今惯于藉两种方式证明真信仰：首先是神圣圣经正典的权威，其次是公教会的传承。这并非因为仅凭圣经正典本身不足以解答一切问题，而是因为许多人按自己的私意解释圣言，采纳各种错误见解，因此必须按照教会信仰的唯一标准来指导对神圣圣经的解释，尤其是在一切公教道理所赖以奠基于其上的那些要点上。\par

[77.] 我们也说过，在教会本身，必须同样重视普世的一致意见和古代的权威，以免我们不是从合一的完整中撕裂而陷入分裂，就是从古代的宗教中堕入异端的新说。我们还说，在这同一的教会古代中，有两件事是那些希望远离异端的人必须非常谨慎而认真地注意的：首先，他们应当查明，在古代，整个天主教会司铎团是否曾以大公会议的权威对所述问题作出过决定；其次，如果出现某个新的问题，而对此没有这样的决定，那么他们应当求助于圣教父们的意见，至少是那些各自在不同时代和地方、保持在共融与信德的合一中、被公认为受认可的导师的意见；并且，凡发现他们同心同德所持守的，就应当毫无疑虑或顾忌地视为教会真实而大公的道理。\par

[78.] 为了不使我们显得是凭自己的权威而非教会的权威妄加断言，我们援引了约三年前在亚洲的厄弗所召开的圣会议为例，当时在巴苏斯和安提奥库斯执政期间，当提出关于权威地确定信仰规则的问题时，唯恐任何亵渎的新奇会潜入，正如在阿里米农所发生的真理的歪曲一样，那里聚集的几乎两百位司铎全体一致赞同这是最符合大公精神、最可信赖、最好的做法，即提出圣教父们的意见，其中一些众所周知是殉道者，一些是精修者，但所有的人都曾是，并且至终仍然是，大公教会的神职人员，以便通过他们的一致决定，适当地、隆重地确认对古代真理应有的敬畏，并谴责亵渎的新奇。这样做之后，那不虔诚的聂斯多略被合法而应得地判定为反对大公的古代传统，反之，真福济利禄则与之相符。而且，为了不使此事缺乏可信度，我们记录了教父们的名字和数目（尽管我们忘记了顺序），根据他们一致而统一的判断，既阐明了诉讼程序的庄严预备，也确立了神圣真理的规则。为了增强我们的记忆，在这里再次提及他们并非多余的工作。\par

% structure: p0165 source=h2 target=section reason=relative-heading-rank; base=h2

\section{第三十章}

\Emph{厄弗所大公会议。}\par

[79.] 这些就是其著作——无论是作为法官还是证人——在会议中被诵读的人：亚历山大主教\Person{圣伯多禄}，一位最杰出的导师和最可敬的殉道者；同城主教\Person{圣亚大纳削}，一位最忠信的教师和最卓越的精修者；同样也是同城主教的\Person{圣德奥斐禄}，一位以其信德、生活、知识而闻名的人，他的继承人、可敬的\Person{济利禄}如今装饰着\Place{亚历山大}教会。唯恐会议批准的教义被认为仅限于一个城市和一个省，还加上了\Place{卡帕多细亚}的那些光亮：\Person{圣额我略·纳齐安}，主教兼精修者；\Person{圣巴西略·凯撒勒雅}，主教兼精修者；以及另一位\Person{圣额我略}，\Person{圣额我略·尼撒}，因其信德、品行、正直和智慧，最配作\Person{巴西略}的兄弟。唯恐希腊或东方显得孤立，为了证明西方和拉丁世界也一直持守同样的信仰，在会议中宣读了\Place{罗马}主教、殉道者\Person{圣斐理斯}和\Person{圣犹利}的若干书信。而且，为了不仅头部，而且世界的其他部分也能为会议的决定作证，从南方加上了最可敬的\Place{迦太基}主教兼殉道者\Person{圣西彼廉}，从北方加上了\Place{米兰}主教\Person{圣安博}。\par

[80.] 所有这些，以神圣的十数，在厄弗所被提出作为导师、参议、证人、法官。而那可敬的会议持守他们的教义，追随他们的忠告，相信他们的证言，服从他们的判断，不仓促，不预断，不偏私，就信仰的规则作出了决定。本可以引用更多的古代教父，但并无必要，因为既不宜让时间被众多证人所占据，也无人认为那十位实际上与其余的同僚意见不同。\par

% structure: p0169 source=h2 target=section reason=relative-heading-rank; base=h2

\section{第三十一章}

\Emph{厄弗所教父们在驱除新异、维护古代方面的恒心。}\par

[81.] 在此之后，我们又加上了真福\Person{济利禄}的判词，该判词包含在同一教会会议记录中。因为当迦太基主教\Person{卡普雷奥卢斯}的书信被宣读后，他在信中恳切请求驱除新异、维护古代，济利禄提出并通过了议案，在此插入该议案也许并非不适时：因为他说，在会议记录结尾，\Quote{也请将已宣读的迦太基可敬而极虔诚的主教卡普雷奥卢斯的书信载入会议记录。他的心意显而易见，因为他恳求确认古代信仰的教义，而将新奇的、肆意编造的、不虔敬的教义予以谴责和弃绝。}全体主教喊道：\Quote{这是所有人的话；这是我们所有人所说的，这是我们所有人所愿望的。}所谓\Quote{所有人的话，}是什么意思？所谓\Quote{所有人的愿望，}是什么意思？岂非就是，从古代传承下来的应该被保持，新近编造的应该被鄙弃地拒绝？\par

[82.] 随后，我们表示了对那次大公会议的谦卑与圣德的钦佩，以致尽管司铎人数众多，其中几乎大部分是都主教，如此博学、如此有学问，以致几乎所有的人都能够参与教义讨论，而他们本身因被召集而来这一事实，似乎可能鼓励他们凭自己的权威做出某些决定，然而他们却没有创新任何事物，没有擅自假设任何事物，绝对没有为自己僭取任何事物，而是尽一切可能小心谨慎，只将他们从教父那里领受的传承给后代。而且，他们不仅满意地处理了当前的事务，也为后来的人树立了榜样，使他们也能坚持神圣的古代传承，并谴责亵渎的新奇。\par

[83.] 我们还痛斥了\Person{聂斯脱利}的邪恶狂妄，他夸口自己是第一个也是唯一理解圣经的人，而在他之前解释神圣教诲的所有教师都是无知的，也就是说，全体司铎、全体精修圣人和殉道者都是无知的，其中一些人曾发表过关于天主法律的注释，另一些人则在他们的注释中同意或默许了这些注释。总之，他自信地断言，整个教会至今仍在错误中，且一直处于错误之中，因为在他眼中，教会一直追随并仍在追随无知和误导人的教师。\par

% structure: p0174 source=h2 target=section reason=relative-heading-rank; base=h2

\section{第三十二章。}

\Emph{切莱斯廷和西斯笃，罗马主教，反对新奇的热情。}\par

[84.] 上述内容已经足够，甚至远远超过粉碎和消灭一切亵渎新奇的所需。但为使这样完整的证据无所欠缺，我们在结尾处添加了宗座的双重权威：首先，是\Person{圣西斯笃教宗}的权威，这位可敬的教宗如今装饰着罗马教会；其次，是他的前任，享真福的\Person{切莱斯廷教宗}的权威。我们认为也有必要在此插入这两项权威。\par

\Person{圣西斯笃教宗}在他就\Person{聂斯脱利}之事致\Place{安提约基雅}主教的\WorkTitle{书信}中说：\Quote{因此，因为正如宗徒所说，信德是唯一的——显然，是迄今所获得的信德——让我们相信将要被说的事物，并说出将要被持守的事物。}什么是将要被相信和将要被说的事物？他接着说：\Quote{不得允许新奇有任何自由，因为给古代传承增加任何东西都是不合适的。不要让祖先的清晰信德和信仰被任何浑浊的混合物玷污。}真正宗徒般的情操！他以清晰的修饰语来提升教父的信德；亵渎的新奇他称为浑浊。\par

[85.] \Person{享真福的切莱斯廷教宗}也以类似方式和相同效果表达了自己。因为在他写给高卢司铎的\WorkTitle{书信}中，指责他们与错误串通，因为他们的沉默使他们在古代信德上失职，并允许亵渎的新奇兴起，他说：\Quote{我们应受责备，如果我们以沉默助长错误。因此，斥责这些人。限制他们讲道的自由。}但此处有人可能会疑惑他禁止谁的任意讲道自由——是古代传承的宣讲者还是新奇的发明者？让他自己告诉我们；让他自己解决读者的疑惑。因为他接着说：\Quote{如果情况如此（即，情况确如某些人向我对你们的城市和省份所抱怨的那样，因你们有害的掩饰使你们同意某些新奇），如果情况如此，就让新奇停止攻击古代传承。}这就是\Person{享真福的切莱斯廷}的判决，不是让古代传承停止推翻新奇，而是让新奇停止攻击古代传承。\par

% structure: p0179 source=h2 target=section reason=relative-heading-rank; base=h2

\section{第三十三章。}

\Emph{天主教会子女应坚持祖先的信德并为它而死。}\par

[86.] 凡是否定这些宗徒的及公教会的决定的人，首先必然侮辱圣策肋定的纪念，他下令新意不得攻击古传；其次，他蔑视圣西斯笃的决定，其判词是：\Quote{不得允许新意的自由，因为不适合对古传作任何添加；}此外，他谴责真福济利禄的决定，济利禄高度赞扬可敬的卡普雷奥卢斯的热情，因为他希望信德的古老教义得到确认，而新奇的发明被谴责；更进一步，他践踏厄弗所大公会议，即几乎整个东方的圣主教们的决定，他们在天主的指引下颁布，后代只应相信圣教父们在基督内一致同意的神圣古传所持守的，他们异口同声、大声宣告，这些是众人的话，这是众人的愿望，这是众人的判词，即正如在聂斯多略之前几乎所有异端者因蔑视古传、支持新奇而被谴责，同样，新奇的作者和古传的攻击者聂斯多略也应被谴责。凡不赞同他们由神圣的、天上恩宠的恩赐所激发的一致决定的人，必然认为聂斯多略的亵渎被不公义地定罪；最后，他蔑视整个基督的教会及其导师、宗徒和先知，特别是真福宗徒保禄，视之为卑贱和无价值；他蔑视教会，因为教会从未失忠于保存一次交付给她的信德的责任；他蔑视圣保禄，他写道：\Quote{弟茂德，要保管所受的寄托，躲避亵渎的新奇言论；}\BibleRef{弟前 6:20} 并且又说：\Quote{若有人给你们宣讲与你们所接受的不同，当受诅咒。}\BibleRef{迦 1:9} 但如果既不能违反训令也不能违反教会的法令——这些法令根据普世性和古传的神圣一致，所有异端者，最后还有贝拉齐、塞勒斯提和聂斯多略，已被正确和应得地谴责——那么肯定所有渴望证明自己是慈母教会的真子女的公教徒，都应从此以后持守圣教父们的神圣信仰，与之结合，死于其中；至于亵渎之人的亵渎新奇——要憎恶它们、厌恶它们、反对它们、不给它们任何宽容。\par

[87.] 这些在前两篇\WorkTitle{备忘录}中更为详尽处理的事项，我现在以总结的方式更简要地汇集在一起，以便我的记忆——为了帮助它我撰写了这些——一方面通过频繁参阅得到更新，另一方面避免因冗长而疲劳。\par

\SourceRecord{Translated by C.A. Heurtley.；Nicene and Post-Nicene Fathers, Second Series；Vol. 11.；Edited by Philip Schaff and Henry Wace.；Buffalo, NY: Christian Literature Publishing Co.,；1894.；<http://www.newadvent.org/fathers/3506.htm>.}
